The journey culminating in this thesis, *Bridging the Agentic Gap*, has been anything but a straight line. If my publication record leading up to this point appears somewhat fragmented, it is a direct reflection of my genuine, unbridled excitement for the field of Artificial Intelligence. I have always been drawn to diverse ideas and collaborations, and while this led me down varied research paths, at the given time and context, pursuing every single one of those threads made perfect sense. Ultimately, it was this wide-ranging curiosity that allowed me to synthesize the dual paradigms presented in this work: the biologically-inspired structural approach of neuroscience and the scalable realm of deep generative world models.

This journey truly began when my original plan to pursue a full-time Master's degree was wonderfully derailed. Prof. Dr. Thilo Stadelmann offered me the opportunity to study part-time while working in his research lab. That decision changed everything. I stayed on through my subsequent studies and into this PhD, always feeling incredibly fortunate for the immense freedom and decision-making power I was given to explore the topics I was most passionate about. Thilo, thank you for your mentorship, for sponsoring my initial fellowship, and for the countless highlights at the ZHAW Centre for Artificial Intelligence. It has been a pleasure being a part of this journey, from the centre's founding to its growth into seven research groups. 

I am equally indebted to the head of my committee, Prof. Dr. Benjamin Grewe, who so warmly welcomed me into his lab. The opportunity to work at the intersection of neuroscience and AI under your guidance was instrumental in shaping the foundational concepts of this thesis. I would also like to express my sincere gratitude to Prof. Dr. Giacomo Indiveri for serving on my committee, as well as to Dr. Jan Milan Deriu, Dr. Jorge Peña Querlata, and Prof. Dr. Christoph von der Malsburg, all of whom provided vital support, insights, and encouragement during different phases of my research.

My research was also made possible through the generous support of the DIZH (Digitalization Initiative of the Zurich Higher Education Institutions), and I am deeply grateful to the Canton of Zurich for awarding me the fellowship that funded the latter stages of this work.

Finally, the most profound thanks belong to my family. Over the course of this PhD, our family beautifully grew from three to five. To my wonderful wife, Dona Belle: thank you for your endless patience, your unwavering support, and for keeping everything together while I was lost in research. To my oldest son, Leeroy, and my daughters, Reyna and Elody: you are my greatest achievements. Thank you for reminding me every day of the real world outside of deep learning models. 
